\documentclass[12pt,a4paper]{exam}

% --------------------
% Packages
% --------------------
\usepackage[utf8]{inputenc}
\usepackage[T1]{fontenc}
\usepackage[french]{babel}
\usepackage{amsmath,amssymb}
\usepackage{geometry}
\usepackage{pgfplots}

\pgfplotsset{compat=1.18}
\geometry{margin=2cm}

% --------------------
% En-tête / pied (exam)
% --------------------
\header{}{DS de mathématiques}{2\up{nde} générale}
\footer{}{\thepage}{\it{G.M.}}

% --------------------
\begin{document}

\begin{center}
    \Large\textbf{Devoir surveillé\\Fonctions affines et vecteurs}
\end{center}

\vspace{0.3cm}

\noindent
\textbf{Durée :} 1h \hfill
\textbf{Calculatrice :} \underline{NON} autorisée

\vspace{0.5cm}

% ======================================================
\noindent
\textbf{Exercice 1 : Vecteurs et parallélogramme (\ \ \ \  / 5 points )}
\vspace{5pt}

On considère la figure suivante :

\begin{center}
    \begin{tikzpicture}[scale=0.7]
        
    % Grille
    \draw[step=1cm, gray, very thin] (-8,-6) grid (14,6);

    % Points
    \coordinate (R) at (-1,-2);
    \coordinate (S) at (5,-2);
    \coordinate (T) at (7,2);
    \coordinate (U) at (1,2);

    % Centre (intersection des diagonales)
    \coordinate (O) at (3,0);

    % Parallélogramme
    \draw[thick] (R) -- (S) -- (T) -- (U) -- cycle;

    % Diagonales
    \draw[dashed] (R) -- (T);
    \draw[dashed] (S) -- (U);

    % Points
    \fill (R) circle (2pt) node[below left] {$R$};
    \fill (S) circle (2pt) node[below right] {$S$};
    \fill (T) circle (2pt) node[above right] {$T$};
    \fill (U) circle (2pt) node[above left] {$U$};
    \fill (O) circle (2pt) node[above] {$O$};

    \end{tikzpicture}
\end{center}

\begin{questions}

\question[1,5] \textbf{Sur le graphique}, placer le point $F$, image du point $R$ par la translation de 
vecteur $\overrightarrow{TU}$.

\question[1,5] \textbf{Sur le graphique}, placer le point $E$, image du point $F$ par la translation de 
vecteur $\overrightarrow{RU}$.

\question[2] À l'aide des propriétés vectorielles, démontrer que $RSUE$ est un parallélogramme.

\end{questions}

\vspace{0.5cm}

% ======================================================
\noindent
\textbf{Exercice 2 : Étude de fonctions (\ \ \ \  / 10 points )}

\vspace{5pt}
On considère la fonction $f$ définie sur $\mathbb{R}$ dont la représentation graphique est donnée ci-dessous.

\begin{center}
\begin{tikzpicture}
\begin{axis}[
    axis lines = middle,
    xmin = -4, xmax = 4,
    ymin = -2, ymax = 5,
    xlabel = {$x$},
    xtick distance=1,
    ytick distance=1,
    ylabel = {$y$},
    grid = both,
    width=14cm,
    height=8cm
]

\addplot[thick,domain=-10:10] {-1/2*x + 2};
\node at (axis cs:2.5,2.5) {$\mathcal{C}_f$};

\end{axis}
\end{tikzpicture}

\end{center}

\begin{questions}

% -----------------------------
\question[1] Quelle est la nature de la fonction $f$ ?

% -----------------------------
\question[2] \underline{En justifiant \textbf{soigneusement}}, trouver l'expression de $f(x)$.

% -----------------------------
\question[2] Résoudre \underline{graphiquement} l’inéquation :
\[
f(x) \geq 3
\]

% -----------------------------
\question On considère la fonction $g$ définie par :
\[
g(x) = 2x + 2
\]

\begin{parts}
\part[2] \textbf{Sur le graphique}, tracer la représentation graphique de la fonction $g$.
\part[1] Déterminer graphiquement les coordonnées du point d’intersection des 
droites représentant $f$ et $g$.
\end{parts}

\vspace{5pt}

% -----------------------------
\question[2] Résoudre \underline{graphiquement} l’inéquation :
\[
f(x) < g(x)
\]

\end{questions}

\vspace{10pt}

% ======================================================
\noindent
\textbf{Exercice 3 : Vecteurs et colinéarité (\ \ \ \  / 5 points )}

\vspace{5pt}

On considère le graphique suivant : 

\begin{center}
    \begin{tikzpicture}[scale=0.7]

    % Grille
    \draw[step=1cm, gray, very thin] (-8,-4) grid (8,4);

    % Origine
    \coordinate (O) at (0,0);

    % Vecteurs
    \draw[->, ultra thick, black] (-4,2) -- (-5,4) node[midway, above right] {$\overrightarrow{u}$};
    \draw[->, ultra thick, black] (4,-3) -- (6,-2) node[midway, above] {$\overrightarrow{a}$};
    \draw[<-, ultra thick, black]  (-1,2) -- (3,4) node[midway, above] {$\overrightarrow{w}$};

    % Point d'origine
    % \fill (O) circle (2pt);

    \end{tikzpicture}
\end{center}

\begin{questions}
    \question[3] \underline{\textbf{Sur le graphique}}, tracer le vecteur :
    \[2\overrightarrow{w}
    - 2\overrightarrow{u}\]

    \question[2] Quels vecteurs semblent être colinéaires ? \textbf{Justifier en revenant 
    à la définition.}
\end{questions}

\vspace{20pt}


% ======================================================
\noindent
\textbf{Exercice 4 (BONUS) : Vers la 1\up{ère} (\ \ \ \  / 2 points )}

\vspace{5pt}

\underline{Point cours :}
\vspace{5pt}

Pour $h \neq 0$, le taux d'accroissement de la fonction $f$ entre $a$ et $a + h$ est donné 
par la formule :
\[
\dfrac{f(a+h)-f(a)}{h}
\]
On considère la fonction carré $f(x) = x^2$.

\begin{questions}
    \question Calculer $f(a+h)$ en fonction de $a$ et $h$.

    \question Montrer que le taux d'accroissement vaut $2a + h$.
\end{questions}

\vspace{20pt}

\begin{center}
    \textit{Fin du devoir}
\end{center}

\end{document}
